\section{Model and main results}\label{sec:overview}

\subsection{The observable experiment}

Let $\mu$ be a known probability measure on $(0,1/2]$.
Each fresh unit independently samples a hidden magnitude $D\sim\mu$
and a hidden sign $S\mid D=d\sim\operatorname{Ber}(q(d))$.
Conditional on $(D,S)$, its observations are iid Bernoulli with mean
$1/2+(2S-1)D$. Different units have independent latent states and
observation tapes. The unknown measurable allocation belongs to
\[
 \mathcal Q_\rho=
 \{q\in L^\infty(\mu):0\le q\le1,\ \theta(q):=\textstyle\int q\,d\mu\le\rho\},
 \qquad 0<\rho\le1/4.
\]
The target is the positive-sign proportion $\theta(q)$.

An algorithm knows $\mu,\rho,B$, but observes neither $D$ nor $S$.
Using an independent random seed and the observed history, it may request
an observation from a fresh or previously observed unit, interleave
units, stop, and return a measurable real estimate. Every requested
observation costs one, including the first on a unit. Units created but
never observed have no statistical effect and may be sampled only when
their first observation is requested. If $T_A$ is total work, define
\begin{equation}\label{eq:model-risk}
 R_{\mu,\rho}(B)=
 \inf_{\substack{A:\ \sup_{q\in\mathcal Q_\rho}\E_qT_A\le B}}
 \ \sup_{q\in\mathcal Q_\rho}
       \{\E_q(\widehat\theta_A-\theta(q))^2\}^{1/2},
 \qquad B>0.
\end{equation}
This is a uniform expected-budget criterion. Stopping may be unbounded;
finite expected work gives almost-sure finite termination under every
admissible population. Estimates may be clipped to $[0,\rho]$ without
increasing risk. In particular, the constant $\rho/2$ has risk at most
$\rho/2$, and randomization permits meaningful procedures even for $B<1$.

\subsection{An entropy--tail selection law}

For a submeasure $\sigma\le\mu$, let $\Gamma_\sigma$ be its homogeneous
dyadic inverse-square cost entropy. With bands
$I_j=(2^{-j-2},2^{-j-1}]$, put
$a_j=\int_{I_j}d^{-2}\,\sigma(dd)$ and $A=\sum_j a_j$.
When $0<A<\infty$,
\[
 \Gamma_\sigma=A[1+\Ent((a_j/A)_j)].
\]
The finite-fragment definition in Section~\ref{sec:additive} also covers
infinite costs and entropies; $\Gamma_0=0$. All logarithms are natural.
Our main result is the following uniform comparison:
\begin{equation}\label{eq:overview-main}
 \boxed{\qquad R_{\mu,\rho}(B)\asymp\mathcal K_{\mu,\rho}(B),\qquad}
\end{equation}
where
\begin{equation}\label{eq:overview-selection}
 \mathcal K_{\mu,\rho}(B)=
 \min\left\{\rho,\inf_{E\ {\rm measurable}}
 \left[\sqrt{\frac{\rho\Gamma_{\mu|E}}{B}}
 +\int_{E^c}\min\left\{1,\frac1{d\sqrt B}\right\}\mu(dd)\right]\right\}.
\end{equation}
Theorem~\ref{thm:entropy-tail-law} states the result formally.
Both comparison constants are universal, simultaneously for every
probability law on $(0,1/2]$, every $0<\rho\le1/4$, and every $B>0$.
No density, regular variation, positive lower support bound, finite
inverse moment, or finite full-law entropy is assumed.

The selected set pays rare-event cost entropy. Its complement pays an
ordinary inverse-gap cost for $d>B^{-1/2}$ and only its probability mass
below that scale. Equivalently, one minimizes
\[
 \zeta((0,1/2])+
 \frac{\sqrt{\rho\Gamma_\sigma}+\int d^{-1}\tau(dd)}{\sqrt B}
 \quad\text{over }\sigma+\tau+\zeta=\mu,
\]
and caps by $\rho$. Hard measurable assignments suffice; the optimizer
need not be an interval. The split is analytical: the estimator never
observes a unit's membership in a selected component.

\subsection{Proof structure}

The upper construction races a likelihood bank against an ordinary
capped classifier on independent observation streams of the same latent
unit. Its null likelihood uses the full law $\mu$. A mixed-work bound
and finite Poisson multilevel allocation give the three-piece upper
bound in Section~\ref{sec:additive}.

The lower proof in Section~\ref{sec:matching-lower} begins with an optimal
finite-band selection. Its ordinary and unresolved pieces are controlled
by a variational null-work modulus. Optimality constrains the gaps of
the retained bands. Grouping comparable costs and trimming a controlled
number of the smallest gaps produces annular components with admissible
conditional rarities and active mass-saturation constraints. Allocating
rarity proportionally to retained entropy makes their lower bounds add
at every budget. Finite-prefix approximation covers laws accumulating
at zero.

All specialized ingredients used in these arguments are proved in this
manuscript. The appendices include the ordinary classifier, weighted
likelihood bank, homogeneous work sum, multilevel allocation, weighted
sign entropy, scale-sensitive gate, null-work transfer, and additive
localization. Standard concentration, minimax, and convex-duality tools
are cited where used. Section~\ref{sec:related-work} distinguishes these
ingredients and their antecedents from the contribution of the joint
selection law.

\subsection{Quantitative stopping certificates}

A complementary representation retains more quantitative information.
The robust cap-$N$ stopping-flow program $P_N(B)$ has exact attained dual
\[
 P_N(B)^2=
 \max_{q_L,q_H,q_v,q_w\in\mathcal Q_\rho}
 \frac{[\theta(q_H)-\theta(q_L)]_+^2}
 {4+B K_N(q_L,q_H;q_v,q_w)}.
\]
Here $K_N$ is optimal stopped signed-likelihood information relative to
$q_v$ per expected observation under $q_w$.
The four priors have separate endpoint, variance, and work roles.
Theorem~\ref{thm:exact-dual} gives the identity and
Proposition~\ref{prop:dual-bellman} gives its Bellman recursion.

Reference absorption and affinity testing give the explicit comparison
\begin{equation}\label{eq:overview-program}
 R_{\mu,\rho}(B)\le P_{N_*}(B)\le16R_{\mu,\rho}(B),
 \qquad
 N_*=\left\lceil64\max\{B,\rho B^2\}\right\rceil.
\end{equation}
The factor sixteen is a program guarantee, distinct from the conservative
constants in \eqref{eq:overview-main}.

A strengthened lower experiment reveals a unit's latent state at its cap.
Its dual value $G_N$ obeys
\[
 0<G_N\uparrow P_\infty,\qquad
 P_N\downarrow P_\infty,\qquad
 G_N/8\le R_{\mu,\rho}(B)\le P_N.
\]
Thus a verified feasible upper policy $U$ and lower dual value $L$ with
$U\le2L$ certify a factor-sixteen approximation. For every fixed input,
sufficiently accurate program solutions eventually meet that criterion.
No uniform convergence rate or bit-complexity claim is made.
Finite-support inputs give genuine finite conic programs.
The eight rational examples in Appendix~\ref{sec:computation} pass at
caps $1$, $4$, or $8$; independent exact checkers verify their feasibility.

The fixed-cap program also has an exact identification threshold:
for $m$ distinct positive magnitudes its infinite-budget bias floor
vanishes exactly when $N\ge2m-1$. Infinite-support laws have a positive
floor at every fixed finite cap
(Theorem~\ref{thm:cap-identifiability}). This is compatible with the
unrestricted problem, whose cap may increase.

\subsection{Reading guide}

The main distributional proof is in
Sections~\ref{sec:additive}--\ref{sec:matching-lower}.
The supporting material is organized as follows.

\begin{center}
\small
\begin{tabularx}{\textwidth}{@{}p{0.48\textwidth}X@{}}
\toprule
Proof component & Location\\
\midrule
Observable classifiers and multilevel allocation
 & Appendix~\ref{sec:uf-foundations}\\
Weighted sign testing, scale gate, annular saturation
 & Appendix~\ref{sec:lower-foundations}\\
Common hard pairs, adaptive work accounting, additive localization
 & Appendix~\ref{sec:transfer-foundations}\\
Stopping-flow representation and observable realization
 & Appendix~\ref{sec:flow-foundations}\\
Exact robust dual and adaptive comparison
 & Appendices~\ref{sec:exact-dual}--\ref{sec:midpoint-dual}\\
Convergent finite certificates and fixed-cap identification
 & Appendices~\ref{sec:genie}--\ref{sec:identifiability}\\
Exact numerical certificates and reproduction
 & Appendix~\ref{sec:computation}\\
\bottomrule
\end{tabularx}
\end{center}
